\documentclass[11pt,a4paper]{article}
\usepackage[T1]{fontenc}
\usepackage[utf8]{inputenc}
\usepackage{lmodern,amsmath,amssymb,amsthm,mathtools}
\usepackage[margin=25mm]{geometry}
\usepackage{booktabs,longtable,array,enumitem,microtype}
\usepackage[hidelinks]{hyperref}
\setlength{\parindent}{0pt}
\setlength{\parskip}{5pt plus 1pt}
\setlist{leftmargin=1.7em,itemsep=3pt,topsep=4pt}
\newtheorem{theorem}{Theorem}[section]
\newtheorem{proposition}[theorem]{Proposition}
\newtheorem{lemma}[theorem]{Lemma}
\newtheorem{corollary}[theorem]{Corollary}
\theoremstyle{definition}
\newtheorem{definition}[theorem]{Definition}
\newtheorem{example}[theorem]{Example}
\newtheorem{remark}[theorem]{Remark}
\newcommand{\Store}{\operatorname{Store}}
\newcommand{\Obs}{\operatorname{Obs}}
\newcommand{\Law}{\operatorname{Law}}
\newcommand{\leaf}{\operatorname{leaf}}
\newcommand{\Cost}{\mathsf{Cost}}
\newcommand{\Ret}{\mathsf{ret}}
\newcommand{\Bind}{\mathsf{Bind}}
\newcommand{\Batch}{\mathsf{batch}}
\newcommand{\Op}{\mathsf{op}}
\newcommand{\Invoke}{\mathsf{invoke}}
\newcommand{\Code}{\mathsf{Code}}
\newcommand{\Ready}{\mathsf{Ready}}
\newcommand{\Uninit}{\mathsf{Uninitialized}}
\newcommand{\Dfin}{\mathcal D_{\mathrm{fin}}}
\newcommand{\den}[1]{\llbracket #1\rrbracket}
\title{Contextual Refinement and Controlled Recovery\\[5pt]
\large Finite stateful replacement through retained probabilistic cells}
\date{1 November 2026}
\begin{document}
\maketitle

\begin{abstract}
A finite family of local probabilistic simulation matrices is sufficient
for replacement in every well-typed closing context of the frozen O4
language, including stored code aliases, repeated interaction with private
random state, adaptive resource changes, disjoint batches and ranked
callbacks. The observation retains public causal trees, public leaf data,
actual capability identities and kinds, and an externally ordered cost
vector. Private stores remain in the semantics throughout composition.

The proof is constructive. Coupled method laws lift to paired raw
histories and genuine O1 native endpoint squares over a state relation.
The common history object carries the reference cost; its reference leg
is cost-exact and its candidate leg is cost-lowering. O1 occurrence
products and tree grafting then implement the context induction. This
construction also passes through the frozen O3 finite-law and QBS
representations. It does not assert a new strict algebra of chosen
coupling witnesses, full abstraction, or all-path guarantees from
probability laws.

A bounded migration example proves a separate residual recovery theorem.
Authorized rollback restores the functional interface without replenishing
spent request credits, consumes recovery authority, and retains a durable
audit effect. All functional continuations agree after explicit disposal
of unused baseline authority; an audit-authorized observer distinguishes
the residual states. The recovery prefix itself remains observable and
costly. The source is intentionally uncompiled, as authorized.
\end{abstract}
\tableofcontents

\section{The result and its frozen dependencies}\label{sec:scope}
The O5 contract~\cite{Contract} requires open replacement, not agreement
for selected closed tests. We keep the finite CF-0.1 model of O3~\cite{O3}
and the typed syntax, private-store extension and ranked code discipline
of O4~\cite{O4}. In particular, an ordinary code value has a public tag
and public data captures; it does not secretly own a capability. Copies
of a token name are aliases, not new owners. Procedures need their actual
input resources at every invocation.

Write $P\preceq Q$ with the \emph{candidate first}: the candidate $P$
preserves public behavior and is no more costly than the reference $Q$.
The order on costs is ordinary coordinatewise $\le$. This is the reverse
of the direction of a cost-lowering native arrow's source-to-target
cost inequality when one writes the source as the reference. The
construction below fixes this orientation explicitly.

\begin{theorem}[The O5 result]\label{thm:main}
Fix the interface and observation regime of Section~\ref{sec:regime}.
Suppose two finite method implementations have the local relational
certificates of Section~\ref{sec:criterion}. Then:
\begin{enumerate}[label=\textup{(\roman*)}]
\item every finite open O4 term preserves the lifted private-state
relation and the candidate-to-reference cost order;
\item every admitted closing context satisfies
$\Obs(C[P])\sqsubseteq\Obs(C[Q])$;
\item each certificate has a realization as a common paired-history
object with native endpoint-square legs in the normalized-content O1
model, and these witnesses compose by O1 occurrence accumulation;
\item the constructed finite certificates induce the corresponding
couplings in the O3 finite-law and QBS execution models;
\item the controlled recovery protocol in Section~\ref{sec:recovery}
satisfies the stated universally quantified residual guarantee, but not
unrestricted audit equivalence or equality with a no-operation prefix.
\end{enumerate}
\end{theorem}
These are sufficient soundness results. No complete characterization of
contextual refinement, mechanized proof, physical recovery guarantee,
unbounded recursion theorem or general concurrent scheduler theorem is
claimed. The current broader literature already proves higher-order
probabilistic contextual refinement in much richer languages; the precise
project-specific increment is explained in Section~\ref{sec:prior}.

\section{Interfaces contexts and observations}\label{sec:regime}

\subsection{The experiment fixed before comparison}
Fix finite public data sorts and tables, a finite pool of actual token
identities, public token kinds, a finite method signature $\mathcal I$,
and a finite common library of ranked procedure templates with ceiling
$R$. Each template captures only finite public data, including already
constructed lower-ranked code values, and invokes only smaller-ranked
templates. Fix all client primitives and their typed resource transitions.
Replacement occurs through declared abstract method/command slots, not
by giving the client different inspectable procedure tags or quoted
bodies. Supplied bodies satisfy $\varepsilon_{\mathrm{impl}}\subseteq
\varepsilon_{\mathrm{decl}}$ and O4's finite acyclic plugging and rank
conditions, including the absence of a template-to-slot-to-same-template
cycle.
A method implementation may be a primitive finite table or a finite
unfolded O4 body at this interface. A body accepting a callback ranges
over the fixed public finite code menu; its certificate must be proved
at every admitted public argument, including that code value. It cannot
assume a theorem about all future clients.

At a public configuration $a$ of interface $E$, ownership is the actual
bundle $\Delta_E(a)$. An outcome may have several ordered leaf
occurrences, with a different public configuration and capability kind
at each. Output token sets must be disjoint subsets of the input token
set, even on zero-weight raw branches. Bind runs the corresponding
continuation at each actual occurrence. Batch starts from an explicit
partition of input ownership; a reserve is an explicit child. These are
the frozen resource rules, not additional restrictions selected to make
this refinement true.

For implementation $i\in\{P,Q\}$, a state-bearing interface is
\begin{equation}\label{eq:expanded}
 \widehat E_i=\coprod_{a\in|E|}\Store_i(\Delta_E(a)),
 \qquad
 \Store_i(\Delta)=\prod_{r\in\operatorname{dom}\Delta}
                         H^i_{r,\Delta(r)}.
\end{equation}
Non-stateful tokens have a one-element private carrier. The public client
can inspect the public configuration $a$, not its $H$ entries. Pure maps carry entries with their
unchanged owned tokens; only declared operations can inspect or change
those entries. A primitive that accesses several entries must own the
corresponding tokens. There is no unrepresented mutable global state
shared behind two disjoint owners.

\subsection{Public content normalization}
Fix, as part of the signature, a common alphabet of public event
constructors and a set of private unary management constructors. Each
implementation's tree map
\[
 \phi_i:T_iX\longrightarrow TX
\]
preserves leaves and public nodes with their ordered children and sends
a private unary node $h(\theta)$ to $\phi_i(\theta)$. Public structural
$\mathsf{Fork}_n$ nodes are preserved. No public event is relabelled
private after seeing the proposed pair. One may choose all events public,
in which case this map is the identity.

\begin{lemma}[The fixed content comparison]\label{lem:normalization}
The maps $\phi_i$ are cartesian natural, leaf-preserving tree-monad morphisms.
In particular,
\[
 \phi_i(\theta[\theta_1,\ldots,\theta_n])
 =\phi_i(\theta)[\phi_i(\theta_1),\ldots,\phi_i(\theta_n)].
\]
\end{lemma}
\begin{proof}
Induct on $\theta$. A leaf selects the corresponding inserted tree.
A public node preserves its constructor and applies the induction
hypothesis independently to its ordered children. A private unary node
is removed, leaving exactly the child's induction equation. This also
shows that the ordered leaf positions are preserved. Relabeling leaves
commutes with these three clauses, proving naturality; the leaf clause
preserves the unit and the displayed equation preserves multiplication.
For cartesianness, given a leaf map $f:X\to Y$, an original tree
$\rho\in T_iY$ and $\tau\in TX$ with
$Tf(\tau)=\phi_i(\rho)$, the preserved ordered leaf bijection assigns
uniquely to each leaf of $\rho$ the corresponding $X$ label of $\tau$.
This is the unique lift in $T_iX$; hence the naturality square is a
pullback. Deleting an arbitrary branching node has no such proof and is
not an admitted erasure here.
\end{proof}
Costs are not erased by $\phi_i$. Neither is the private store erased at
this stage. The latter is a final observation projection, not a
state-free interpretation of an intermediate operation.

\subsection{The full admitted context class}\label{sec:contexts}
Let $\mathcal C_{\mathcal I,\mathfrak R}$ consist of \emph{all} finite
well-typed closing O4 command contexts with the above fixed interface,
signature, code library and observation regime $\mathfrak R$. An open
method parameter can occur repeatedly. Closing instantiates its declared
public arguments, supplies the actual owned input resources, plugs finite
acyclic implementations, and supplies an admitted initial state law.
An initial law describes the boundary state with accumulated execution
cost zero there. If initialization is part of the observed experiment,
its factory is a declared operation subject to the same local criterion,
with its own events and costs; L1 does not erase that prefix.
The client syntax includes
\[
 \Ret(f),\quad\Op_e(v),\quad\Bind(t,b.u_b),\quad
 \mathsf{case}(v;\vec t),\quad\Batch_\chi(\vec t),\quad
 X[\sigma],\quad\Invoke(p).
\]
All O4 pure data/evidence operations, dependent reindexing, public code
inspection, data framing, and stored aliases remain available. The class
is not a finite list of test clients. At fixed rank its finite derivations
can still be arbitrarily large. A particular service counter further
bounds its effective calls; the general theorem also permits finite
straight-line uses of operations without such counters.

The class permits adaptive queries, adaptive stopping with the remaining
owned token returned, branches with different continuation resources,
repeated calls through stored code aliases, and every lower-ranked
callback typed by O4. It does not admit self-recursive reentry, fresh
token creation, implicit disposal, raw-witness inspection, private-store
projection or effect exchanges absent from O4. A returned native code
tag remains publicly distinguishable. No blanket code opacity is used.

\subsection{Observation and its order}\label{sec:observations}
For a completed run $h$ of a closed term, retain cost
$c_i(h)\in\mathbb N^d$ and its tree $\theta_i(h)$ with full semantic
leaf configurations. Let $\operatorname{pub}_F:\widehat F_i\to|F|$
forget private entries but retain all public data, token identities and
kinds. The closing observation is the finite law
\begin{equation}\label{eq:observation}
 \Obs(t)=\sum_hp_i(h)\,
 \delta_{\left(c_i(h),T\operatorname{pub}_F\phi_i(\theta_i(h))\right)}.
\end{equation}
It retains public causal order and ordered leaf positions, not merely the
terminal value marginal. If a public intermediate response must appear
in this external observation, its method's declared event label includes
that response; otherwise it is available to its continuation without
being automatically added to the event label. Either convention is fixed
by the signature. The examples below explicitly use response-labelled
events.

For finite laws $\mu,\nu$ on $\mathbb N^d\times T|F|$, set
\begin{equation}\label{eq:obsorder}
 \mu\sqsubseteq\nu
 \quad\Longleftrightarrow\quad
 \exists\gamma\in\Dfin((\mathbb N^d\times T|F|)^2):
 \begin{cases}
  (\pi_1)_*\gamma=\mu,\quad(\pi_2)_*\gamma=\nu,\\
  \gamma((c,t),(c',t'))>0\ \Longrightarrow\ t=t',\ c\le c'.
 \end{cases}
\end{equation}
The equality regime strengthens $c\le c'$ to $c=c'$. An exact-cost
coupling then yields equality of the joint observations. In the ordered
regime the public marginal is equal and, for any bounded increasing cost
function $g$ with the public tree fixed,
$\mathbb E_\mu g\le\mathbb E_\nu g$. In particular
$\Pr_\mu[c\le b]\ge\Pr_\nu[c\le b]$.

Cost is external in this order. If a client can read a reported clock or
budget as ordinary data, that value must match \emph{exactly} in the
public relation; lower accumulated cost alone cannot justify a different
branch decision. The law makes no all-raw-path assertion about zero
probabilities. Every admitted command mathematically terminates by its
finite syntax and rank bound; there is no divergent behavior, scheduling
choice or fairness hypothesis in the theorem. The causal tree is not the
O4 phase evaluator's chronological word, and no equality of those words
is inferred.

\section{A finite local simulation criterion}\label{sec:criterion}

\subsection{State relations and their tree lifting}
For each actual token and kind choose a finite relation
$S_{r,k}\subseteq H^P_{r,k}\times H^Q_{r,k}$. Define
\begin{equation}\label{eq:state-relation}
 S_E=\{((a,s),(a,t)):\forall r\in\operatorname{dom}\Delta_E(a),
                     \ (s(r),t(r))\in S_{r,\Delta_E(a)(r)}\}.
\end{equation}
This is a product-shaped \emph{predicate}, not a product probability law.
An initial coupling $\lambda$ can correlate every token's private state
with the others and with public data, as long as it is supported on $S_E$.
The relation restricts to disjoint sub-bundles and reassembles on their
union. A more general relation can replace this sufficient choice only
if those restriction and reassembly properties are separately proved.

Two normalized outcomes $z=(c,\theta)$ and $z'=(c',\theta')$ are
$\mathcal L(S_F)$-related if $c\le c'$ and there is a paired-leaf tree
$\vartheta\in T(S_F)$ with
\begin{equation}\label{eq:tree-lifting}
 T\pi_P(\vartheta)=\theta,\qquad T\pi_Q(\vartheta)=\theta'.
\end{equation}
Equivalently they have the same public constructor shape and their
corresponding actual leaves have the same public configuration and
$S_F$-related private stores. For planar trees this paired tree is unique
once its two projections are fixed. Nullary trees and arbitrary finite
numbers of leaves are included. Related public indices entail equal
result-dependent capability bundles at every matched occurrence.

For each full semantic input $x$ write
\[
 K_i(t)(x)=\sum_{h\in\Omega_i(t,x)}p_i(h)
                 \delta_{(c_i(h),\phi_i\theta_i(h))}.
\]
This retains private leaves. A command certificate
$t_P\preceq_{S_E,S_F}t_Q$ is a finite outcome coupling supported on
$\mathcal L(S_F)$ for \emph{every} pair in $S_E$. Empty relation fibres
have no obligation. No closing observation is taken within this judgment.

\subsection{The independently checkable premises}\label{sec:local}
The local method criterion is the following finite data.
\begin{enumerate}[label=\textup{L\arabic*.}]
\item The intended initial laws $\alpha_P,\alpha_Q$ have a joint coupling
$\lambda$ supported on $S_E$.
\item For every imported method $m$, every admitted common public
argument $v$, and every related full input pair $x\,S_{E_m}\,y$, provide
a nonnegative matrix $G_{m,v,x,y}(z,z')$ with
\begin{equation}\label{eq:local-matrix}
 \sum_{z'}G(z,z')=K_P(m[v])(x)(z),\qquad
 \sum_zG(z,z')=K_Q(m[v])(y)(z'),
 \qquad G(z,z')>0\Rightarrow z\,\mathcal L(S_{F_m})\,z'.
\end{equation}
The target fibre and all resource transitions are the declared ones.
\item Every fixed client primitive has the same property. Identical
primitive names are not sufficient: the actual transition must preserve
$S$. A usual sufficient case uses the same fresh random witness, performs
the same public action, and carries untouched related private entries.
An audit primitive that reveals a forgotten entry must be checked and
may fail.
\item The common pure syntax acts only on public data and carries
unchanged private entries with their owners; the common code library has
identical public tags, captures and templates on both sides. Its bodies
may contain the imported method parameters.
\end{enumerate}
These premises mention method tables, related input states and public
arguments, not arbitrary clients or continuation programs. Fix each
chosen matrix as a function of its actual method/command and related
input pair; the compositional construction does not give it undeclared
access to earlier hidden history. If a method
is a finite compound body, its table can be obtained by the O4 finite
unfolding or by the compositional rules proved below. All sets in a
particular certificate are finite. With rational probabilities, existence
of $G$ is a finite linear feasibility problem after the forbidden pairs
have been set to zero. With arbitrary real weights this is a mathematical
certificate condition, not a claimed decision procedure.

The criterion is stronger than equality of public one-call marginals.
The related private output must still support the next interaction. It
also permits different private carrier cardinalities and nonfunctional
couplings, so it is not merely raw isomorphism in another notation.

\section{The open replacement theorem}\label{sec:replacement}

\begin{theorem}[Open contextual soundness]\label{thm:replacement}
Under L1--L4, for every finite O4 derivation
$\Xi;\Gamma\vdash t:E\Rightarrow F!\varepsilon$, every common public
valuation, and every pair $x\,S_E\,y$, the two interpretations obtained
by plugging $P$ and $Q$ admit a coupling
\[
 K_P(t)(x)\ \mathcal L(S_F)^\dagger\ K_Q(t)(y).
\]
Here the dagger means existence of a finite coupling with the indicated
support, not an additional language operation. Consequently, for every
$C\in\mathcal C_{\mathcal I,\mathfrak R}$ with an L1 initialization,
\begin{equation}\label{eq:replacement}
 \Obs(C[P])\sqsubseteq\Obs(C[Q]).
\end{equation}
If every local comparison is cost-exact, the joint observation laws are
equal. Pure dependent substitution and finite acyclic implementation
plugging preserve these conclusions.
\end{theorem}
\begin{proof}
We prove the stronger open related-input claim simultaneously by template
rank and by the finite typing derivation. At top level use structural
induction; an invocation appeals to a strictly smaller template rank.
All couplings below depend on the actual paired full input, not only on
its public projection.

\emph{Return and pure evaluation.} Equal public valuations evaluate the
same finite pure term to the same data, evidence and code tag/captures.
The pure interface map preserves actual ownership and carries each
private entry unchanged. Its paired result lies in $S_F$. The single
paired leaf with probability one and equal zero cost is the certificate.
The same argument proves preservation under public data framing,
projections, dependent-pair operations and pure reindexing.

\emph{Primitive and imported method.} Use L2 or L3 at the common public
argument and the actual related input pair. Full resulting private
stores remain in the outcomes. This is the implementation-specific base
case, not an appeal to a desired contextual conclusion.

\emph{Case.} The selector is a pure public value and agrees on the two
sides. Both select the same typed branch. Restrict to its actual
public/dependent fibre and use the induction hypothesis. No selector can
inspect a private entry in the frozen pure grammar.

\emph{Bind.} Let $G_0(z,z')$ couple the prefixes. For every positive
pair, write $z=(c,\theta)$ and $z'=(c',\theta')$. Their normalized
public shapes agree and their leaves are paired at the same ordered
occurrence addresses. At matched address $j$, let $x_j\,S_F\,y_j$ be
the full leaf pair. Its public component chooses the same continuation
$u_j$, with exactly the same dependent index and resource bundle. By
induction choose a coupling once for each continuation and actual paired
leaf input, writing $G_j^{z,z'}=G_{u_j,x_j,y_j}$. Its arguments include
any public lexical frame explicitly carried at that leaf; it has no
extra dependence on erased prefix history. Define the finite conditional
product
\begin{equation}\label{eq:bind-coupling}
 G_0(z,z')\prod_{j\in\leaf(\theta)}G_j^{z,z'}(w_j,w'_j),
\end{equation}
and push it forward by grafting both trees and adding their respective
costs. Empty products have mass one and add zero cost.

To verify the $P$ marginal, first sum each $w'_j$. Its coupling marginal
is exactly $K_P(u_j)(x_j)$ and depends only on that left input. The paired coupling itself
is indexed only by its actual paired leaf input. Summing $z'$
then gives the prefix law $K_P(t)(x)$. These are precisely the frozen
O3 occurrence-composition probabilities. The $Q$ marginal is symmetric.
Finite Fubini justifies the reorderings. Prefix and child cost inequalities
add. The paired tree is obtained by grafting the $T(S)$ child trees into
the paired prefix tree, so its leaves remain related. Lemma
\ref{lem:normalization} commutes this graft with the fixed content map.

This calculation conditions on the \emph{entire paired prefix outcome}.
Correlations already carried among its private entries or leaves remain
in that mixing distribution; they are not replaced by independent
marginals. Only fresh continuation runs at distinct actual occurrences
are independent conditional on those inputs, exactly as in O3. Equal
leaf values at different positions do not identify the positions or their
new draws. O4's explicit data frame and actual leaf projection handle
persistent lexical values without duplicating the original ownership.

\emph{Batch.} At a fixed related global input pair, the common public
splitter $\chi$ yields equal disjoint child token bundles. Restrict the
paired stores and $S$ to those bundles. The product of the child
inductive couplings, including any reserve child, has the correct
conditional marginals. Graft under the common public $\mathsf{Fork}$
and sum costs. Reassembly of the tokenwise relation proves the leaf
condition. If the global input is random, integrate this conditional
product against its original \emph{joint} input coupling. In particular,
we never factor a correlated initial store law into independent token
laws. The product applies to fresh transitions conditional on the full
paired input, not to the private input states themselves.

\emph{Open calls and plugging.} The actual public parameter substitution
selects the same arguments. An imported method uses L2. A compound
implementation already certified by the induction can be used at that
same interface. O4's typed plugging theorem guarantees that finite
capture-avoiding substitution preserves each actual resource and result
fibre. The allowed assignment is acyclic; substitution creates no
recursive call knot and no owned capture.

\emph{Stored code and ranked invocation.} Equal public code values have
the same inspectable tag and captures, regardless of how often the data
were copied or stored. Invocation selects the same template and
instantiates it by the same pure substitution. Its lower rank supplies
the induction hypothesis at the current related owned stores. Reusing an
alias later invokes at the later store; the preceding step has already
proved that store related. New sampling inside an invocation uses its
fresh coupled transition, whereas an earlier stored random value stays
in the public data or private store. The proof does not replace copying
a sampled value by rerunning a sampler. Every admitted lower-ranked
callback is another such invocation. Unbounded reentry is outside the
syntax, not a missing induction case.

For dependent substitution, evaluate each premise at the substituted
public valuation. O4's substitution theorem reindexes the actual
capability family as well as the data. The same finite coupling serves
at that reindexed pair. This also handles parameter instantiation and
local bind-variable projections. The preceding plugging case proves
stability under finite simultaneous implementation substitution.

Finally mix the open coupling against $\lambda$ from L1. Summing either
input coordinate yields its stipulated initial law. Push the two output
coordinates through $T\operatorname{pub}_F$: related private leaves
become equal public leaves, and the cost inequality is unchanged. This
is exactly~\eqref{eq:obsorder}. With exact local costs, every addition
preserves equality, giving equality of the joint closing laws.
\end{proof}
This proof establishes the quantifier over the whole context class.
Enumerating example clients, even every client below a chosen syntax
size, would not establish it. Conversely, the theorem does not say that
every contextually valid replacement has one of these tokenwise local
simulation certificates.

\section{Realization by the contracted categorical cells}\label{sec:cells}

\subsection{The native cells that must actually be used}
The frozen joint rig is
\[
 \mathcal V=(\mathbb R_{\ge0})_{\mathrm{disc}}\times\Cost,
 \quad(p,c)\oplus(q,d)=(p+q,\min(c,d)),
 \quad(p,c)\otimes(q,d)=(pq,c+d),
\]
where $\Cost=(\mathbb N\cup\{\infty\})^d$ has an arrow $c\to d$
when $c\ge d$ coordinatewise. All execution costs here are finite. In
the O1 mixed construction~\cite{O1}, a retained computation has a finite
augmented witness family $(\Omega,p,c)$ and a tree body $\theta$.
For endpoint maps $u:A\to A'$ and $v:B\to B'$, a native square has
witness functions $r_a$ satisfying
\begin{equation}\label{eq:native}
 p'_{u(a),w'}=\sum_{r_a(w)=w'}p_{a,w},\qquad
 c'_{u(a),w'}\le\min_{r_a(w)=w'}c_{a,w},\qquad
 Tv\,\theta_a(w)=\theta'_{u(a)}(r_a(w)).
\end{equation}
The empty minimum is $\infty$. Cost-exact means equality of costs on
every source witness, not surjectivity onto all target witnesses.
O3 uses exactly these equations, not an independently invented notion of
cell. The concrete content comparison induced by $\phi_i$ changes the
body to $\phi_i\theta$ and leaves raw witnesses, probabilities and costs
unchanged. Lemma~\ref{lem:normalization} is the multiplication condition
needed for this comparison. More explicitly, the ordered-hole
bijection identifies each dependent product in~\eqref{eq:occurrence}
before and after normalization; the witness tuples, weights and costs
are unchanged and grafting commutes. This proves the strong comparison
on this terminal-grade instance directly, without appealing to a
change-of-input theorem requiring an invertible extension map.

\subsection{Lifting a law matrix back to retained histories}
Fix a related full input pair $x\in S_E$ and a local or derived outcome
coupling $G$. Let $w\in\Omega_P$ and $v\in\Omega_Q$ be the original
raw histories, with probabilities $p_w,q_v$, normalized full outcomes
$a(w),b(v)$, and pushforward laws $\mu,\nu$. Define
\begin{equation}\label{eq:raw-lift}
 \Gamma(w,v)=
 \begin{cases}
 G(a(w),b(v))\dfrac{p_w}{\mu(a(w))}\dfrac{q_v}{\nu(b(v))},
       &\mu(a(w))\nu(b(v))>0,\\[6pt]
 0,&\text{otherwise}.
 \end{cases}
\end{equation}
This is a finite conditional-weight formula, not a measure-theoretic
choice of arbitrary disintegrations.

\begin{lemma}[Raw lift and common cost]\label{lem:raw-lift}
The matrix $\Gamma$ has marginals $p,q$. Its positive pairs define a
finite normalized common retained object
$W:S_E\Rightarrow S_F$, with a unique paired-leaf body and cost
$c_W(w,v)=c_Q(v)$. There are native endpoint squares
\begin{equation}\label{eq:span}
 \begin{array}{ccc}
 &W:S_E\Rightarrow S_F&\\[-2pt]
 {}_{(\pi^E_P,\pi^F_P)}\swarrow&&
       \searrow_{(\pi^E_Q,\pi^F_Q)}\\[-2pt]
 H_P:E_P^{\wedge}\Rightarrow F_P^{\wedge}
 &&H_Q:E_Q^{\wedge}\Rightarrow F_Q^{\wedge},
 \end{array}
\end{equation}
after the specified content normalization. The reference leg is
cost-exact; the candidate leg is cost-lowering.
\end{lemma}
\begin{proof}
For $p_w>0$, summing first over raw $v$ with a fixed outcome $b$ turns
$\sum_{b(v)=b}q_v/\nu(b)$ into one. The coupling marginal then gives
\[
 \sum_v\Gamma(w,v)
 =\frac{p_w}{\mu(a(w))}\sum_bG(a(w),b)=p_w.
\]
If $p_w=0$, every summand is zero, so the equation still holds. The
other marginal is symmetric. Hence the positive carrier has total mass
one. Each supported pair has aligned normalized trees and related leaf
stores, so~\eqref{eq:tree-lifting} supplies its paired body.

The history functions are $(w,v)\mapsto w$ and $(w,v)\mapsto v$.
Their probability pushforward equations are exactly the two marginals.
Their body equations are~\eqref{eq:tree-lifting}, with input maps the
relation projections. On every common history the reference target cost
is $c_Q(v)=c_W(w,v)$. The candidate target cost satisfies
$c_P(w)\le c_Q(v)=c_W(w,v)$. Therefore its target cost is at most every
cost in its history fibre and hence at most their minimum. A target with
no preimage has weight zero and satisfies the cost inequality against
$\infty$. Nonnegative finite common costs give the required augmentation
$(\sum\Gamma,\min c_W)\to(1,0)$. Thus all three equations
in~\eqref{eq:native} hold on the actual O1 carrier.
\end{proof}
This is a proof-relevant span of native endpoint squares over a state
relation. It need not be a direct cost-lowering native cell $H_Q\to H_P$: their
private endpoint carriers can differ, and the coupling need not be the
graph of a function. The relation carrier is mathematical proof data,
not a program-owned resource or an operation exposing private state.
The relation projections are mathematical endpoint maps in O1's ambient
all-Set construction; they are not asserted to be executable client pure
maps or authorized physical capability transformations.
If private unary events differ, the squares live after the fixed content
comparison; no full unnormalized body equation is asserted.

\subsection{Composition is the actual occurrence construction}
For a prefix witness $h$ whose tree has addressed holes $j$, O1 uses
\begin{equation}\label{eq:occurrence}
 \Omega_{t;u}(a)=\coprod_{h\in\Omega_t(a)}
                 \prod_{j\in\leaf(\theta_t(h))}\Omega_u(b_j),
 \qquad
 p=p_t(h)\prod_jp_u(k_j),\quad c=c_t(h)+\sum_jc_u(k_j),
\end{equation}
with tree body obtained by grafting. Apply this formula to the paired
prefix object $W_t$ and the paired continuation objects at its actual
$S_F$ leaves. Its histories retain both component histories at every
address. Its cost is the sum of their reference costs, hence the full
reference cost. The two endpoint projections send the paired tuple to
its $P$ and $Q$ tuples. Finite distributivity proves their marginal
equations; tree grafting proves their body equations; addition proves
the exact or lowering cost equations. Thus the bind construction in
Theorem~\ref{thm:replacement} is literally an O1 occurrence composite,
not a callback operational argument with a categorical name appended.
The same product-and-fork calculation gives the separated batch square.

There is no requirement that the raw lift of the final pushed-forward
coupling be the \emph{same chosen witness set} as this composite span.
Both realize the same needed relational statement. The composite span
is a concrete retained certificate with its own histories; preservation
of the theorem is existence of such a certificate, not equality of every
choice of conditional splitting.

\begin{proposition}[Reflexivity and vertical gluing]\label{prop:vertical}
The observation order is a preorder. Local relational certificates
compose when their state relations compose, and the existence judgments
are stable under the horizontal language constructions proved above.
\end{proposition}
\begin{proof}
Reflexivity uses diagonal couplings and equality of private states.
For laws $\mu,\nu,\xi$ with couplings $G(a,b)$ and $D(b,c)$ define
\begin{equation}\label{eq:gluing}
 J(a,b,c)=
 \begin{cases}G(a,b)D(b,c)/\nu(b),&\nu(b)>0,\\0,&\nu(b)=0.
 \end{cases}
\end{equation}
Zero middle mass forces both associated marginal fibres to be zero.
Summing gives the two input couplings and, after marginalizing $b$, a
coupling of $\mu,\xi$. Equal public tree shapes and transitive cost
order give the observation preorder.

For state-indexed local certificates, let
$S^{13}_{r,k}=S^{12}_{r,k};S^{23}_{r,k}$. At any related endpoint input
pair, choose and retain a finite middle private-state witness at every
owned token, giving one middle full input. Use the two local matrices
at those two input pairs, and apply~\eqref{eq:gluing} to their common
middle full-outcome law. Every resulting leaf has a middle state witness
and thus belongs to $S^{13}$. Reassembly works because the states follow
the same disjoint owned tokens. The middle input and outcome witnesses
must not be dropped before this construction.

Iterated vertical gluing can retain all middle variables; its positive
joint masses are products of adjacent coupling masses divided by the
positive intermediate marginals. Rebracketing the finite sums yields
the same outer law when these choices are fixed. Horizontal closure is
the separate occurrence proof already given.
\end{proof}
The safe categorical conclusion is an observation-directed,
proof-irrelevant comparison of existence judgments, together with the
explicit proof-relevant native spans that realize it. We do \emph{not}
assert a strict double category of chosen coupling witnesses, a chosen
horizontal--vertical interchange law, an equipment, or a faithful
encoding of all O1 parallel cells. An arbitrary native hom cell whose
endpoint/body relation does not preserve the selected public interface
is not thereby a sound replacement. A native functional cell satisfying
that interface condition gives a special graph coupling; the full
relation criterion is the additional compatibility needed by O5.

\subsection{Zero probability and the exact all path boundary}
\label{sec:zeros}
The positive carrier in Lemma~\ref{lem:raw-lift} need not cover
zero-weight raw target histories. Thus exact joint laws and even a
cost-exact reference leg do not prove an all-path bound. For example,
add to an otherwise cost-zero candidate a raw branch of probability
zero and cost $100$, with the same public body. Its observed law is
unchanged; its all-raw-path budget $0$ is false.

\begin{proposition}[The additional coverage guarantee]\label{prop:coverage}
If a closing paired-history span has a surjective candidate-history
projection, then an all-raw-path upper cost bound for the reference is
an all-raw-path upper bound for the candidate. A sufficient local
strengthening is candidate coverage of every local raw method and
client-primitive certificate, including zero-weight histories, together
with initial candidate-state coverage and the same support/body/cost
conditions on those added pairs. This strengthening is preserved by the
admitted context constructors.
\end{proposition}
\begin{proof}
Lift any candidate raw history to a common history. Its candidate cost
is at most the common cost, which equals the projected reference cost.
A reference all-path bound therefore bounds the candidate history.
For the local assertion, induct on an arbitrary candidate run. Return,
case and ranked dispatch select the same public branch. At an operation,
local coverage gives a paired reference raw history, even if its weight
is zero. At bind, lift the prefix and then each candidate child history
at the corresponding related leaf; the fixed tree comparison preserves
all its positions. Batch lifts each disjoint child and retains the
reserve. This yields a common composite history and proves surjectivity.
Initial coverage supplies the needed related starting pair.
\end{proof}
The extra zero-mass pairs may be adjoined only when their body and cost
inequalities really hold. Probability normalization alone supplies no
such pairs. The main replacement theorem uses the law/almost-sure regime
and does not assume this optional coverage strengthening.

\section{The finite representation comparison}\label{sec:representation}
\begin{theorem}[Compatibility with the frozen representations]
\label{thm:representation}
Every paired-history certificate constructed above has compatible
finite-law, operational causal/cost and finite-discrete QBS images.
Its two marginal laws are the corresponding interpretations of the two
programs; its support proves the same closing observation order.
\end{theorem}
\begin{proof}
For a finite augmented witness family, the O3 comparison is
\[
 L(H)(a)=\sum_{w\in\Omega(a)}p_w\delta_{(c_w,\theta_w)}.
\]
For the common object, retain in addition the two projected costs and
bodies, yielding the explicit joint pushforward
\[
 \sum_{(w,v)}\Gamma(w,v)
       \delta_{((c_P(w),\phi_P\theta_P(w)),
                 (c_Q(v),\phi_Q\theta_Q(v)))}.
\]
The native marginal equations prove that this is the required coupling.
The paired body's projection equations prove its state relation. The
common object's reference cost is essential to the native realization;
the displayed two-cost law is the relational comparison, not a false
claim that both native legs preserve exact cost.

Under occurrence composition, the product-and-graft formula
\eqref{eq:occurrence} pushes forward to the O3 law composition on
$\Dfin(\mathbb N^d\times T(-))$. Finite Fubini gives the same coupled
formula~\eqref{eq:bind-coupling}; normalization commutes by Lemma
\ref{lem:normalization}. O4's raw-run correspondence identifies these
with the causal/cost projection of its operational evaluations.

Embed the finite carriers and finite probability tables into discrete
QBS. Each matrix, projection, tree constructor and cost operation is a
map on these finite/discrete objects. The same finite sums have the same
marginals in QBS probability, and the represented relation contains their
support. Thus the certificate lives in the O3 execution model
$KX=P(\mathbb N^d\times TX)$, with its specified independent-occurrence
multiplication. The finite reindexing and plugging equations are
preserved by these same maps.
\end{proof}
This is not a theorem choosing measurable couplings from arbitrary
pointwise QBS existence assertions; every coupling family here was a
finite table. Nor is it a theorem about the undecorated Markov monad
$P$ alone: the execution meaning retains the tree and cost. Projecting
private state after \emph{each} method generally breaks composition.
The theorem erases it only when closing, after all interaction. A
state-free per-method quotient would require a separate compatibility
proof and is refuted in general by the next example.

\section{Implementation change and an interaction distinguisher}
\label{sec:examples}

\subsection{A genuinely different cheaper persistent implementation}
Fix a token $r$, a bound $N$, and public kinds $\Uninit$ and
$\Ready_h$ for $0\le h\le N$. The reference $Q$ privately stores
$(b,j)\in\{0,1\}^2$, while the candidate $P$ stores
$b\in\{0,1\}$. Relate
\[
 b\ S_{r,\Ready_h}\ (b,j)\qquad(j\in\{0,1\}).
\]
At $\Uninit$, both private carriers are singletons, related uniquely.
An initialization method on the preexisting $\Uninit$ token draws a
fair $b$ in $P$, and independent fair $b,j$ in $Q$, and retypes that token
$\Ready_N$. It has cost zero and a private unary management event, fixed
as private in this API. Couple the initial raw branches by
$\Pr(b,(b,j))=1/4$. There are two candidate raw histories and four
reference histories; this is not a raw bijection.

For $h>0$, the public call method returns $b$, retains its private store,
retypes the same token to $\Ready_{h-1}$, and emits
$\mathsf{call}(b)$. Give the candidate cost $1$ and the reference cost
$2$. At every related input pair the Dirac pair of transitions is a
local certificate. A matched explicit disposal method, if included,
consumes $r$ on both sides at the same cost. Fixed unrelated client
primitives preserve the relation by carrying the service entry with its
single owner; any primitive that could inspect it must instead pass L3.
No primitive reports the differing external call costs as public data.

\begin{corollary}[Persistent implementation replacement]
\label{cor:persistent}
For every admitted O4 client at this API, including repeated calls through
stored wrapper-code aliases, adaptive stopping, disjoint unrelated work
and ranked callbacks, replacing $Q$ by $P$ preserves its public causal
law and weakly lowers accumulated cost. On paired runs with $m$ effective
calls, the reference cost exceeds the candidate cost by exactly $m$.
\end{corollary}
\begin{proof}
The initialization matrix and each method pair satisfy L1--L4; the local
call cost gap is one and all other paired costs are equal. Apply
Theorem~\ref{thm:replacement}; its coupling preserves the complete
public tree and therefore the number $m$ of public call nodes. Add the
$m$ local gaps.
\end{proof}
The persistent bit is not resampled by repeated invocation. Common public
wrapper code may call the abstract method slot, and that same wrapper can
be copied, inspected and reused. Replacing the private implementation of
that slot is not returning a different publicly inspectable wrapper tag.
Adaptive resource change is real: if a client calls again only after a
zero response, its final token can be $\Ready_{N-1}$ or
$\Ready_{N-2}$ as a result-dependent public output. The relation and the
proof track those two fibres rather than forcing one fixed continuation
budget. The general theorem also covers multi-hole dependent operations;
this small instance is not the definition of the context class.

If the two call costs are both one, the same construction proves exact
joint-law equivalence despite the different private carriers and raw
histories. If the initial stores of this service and other tokens are
correlated, the theorem still applies whenever a joint $S_E$-supported
initial coupling is supplied; its proof does not require independently
initializing every token.

\subsection{Matching coarse endpoints do not support replacement}
Take two services with the same owned-token protocol
$\Ready_2\to\Ready_1\to\Ready_0$, the same public call labels and
unit call costs. The shared service samples a fair bit $B$ once and
returns $B$ on each call. The fresh service draws a fair bit anew on each
call. Both have the same one-call public marginal and the same coarse
resource endpoints. Their two-call laws are
\[
\begin{array}{c|rrrr}
 &00&01&10&11\\\hline
 \text{shared}&1/2&0&0&1/2\\
 \text{fresh}&1/4&1/4&1/4&1/4.
\end{array}
\]
The well-typed client stores its first answer $x$, calls again with the
actual updated $\Ready_1$ ownership, and returns $(x=y)$ together with
$\Ready_0$. Its probability of true is $1$ versus $1/2$. It can call
through two aliases of one common wrapper code; the aliases do not
supply two owners. Hence neither replacement direction holds for
\eqref{eq:obsorder}.

The exact local obstruction occurs after the first call. The retained
shared state fixes the next answer to $B$, while the fresh service's
conditional next-answer law remains fair. No matrix with those marginals
can place all its mass on equal public replies. Forgetting $B$ after the
first call would manufacture a false memoryless simulation. This is why
the relational criterion keeps full private leaves until closing.

\subsection{Public code and clock boundaries}
If two implementations return different native procedure tags, a pure
case returning the tag's Boolean discriminator separates them, even if
invoking the two procedures would have equal output laws. The present
criterion rejects that pair because its public leaves differ. An opaque
adapter may be separately declared, but cannot retroactively remove
these old code-inspection contexts.

Likewise, if a method returns its exact cost as ordinary data, a context
can branch on that value. Then the positive cheaper-service example does
not satisfy the same public interface relation unless the returned values
agree. The theorem orders only the external accumulated annotation and
is not a license to hide a client-readable measurement.

\section{Controlled recovery with a durable residual}\label{sec:recovery}

\subsection{A bounded protocol with explicit recovery authority}
This is one finite deployment/migration stress case, not an expansion of
the later O6 application programme. Fix actual identities $r$ for a
service, $k$ for a one-shot migration authority, and optionally $a$ for
an audit authority. These identities preexist. Fix $N\in\mathbb N$.
The public kinds of $r$ are $A_h$, $\mathsf{Stage}_h$ and $B_h$ for
$0\le h\le N$. The kinds of $k$ are $\mathsf{Fresh}$ and
$\mathsf{Armed}$; those of $a$ are $\mathsf{Audit1}$ and
$\mathsf{Audit0}$. Other carried resources use their fixed declared
interfaces and do not access $r$'s private audit entry.

The service private state is $(\mathsf{phase},x,d)$, with Boolean
physical payload $x$ and durable audit bit $d$. Its logical value is $x$
in phase $A$ and $1-x$ in phases $\mathsf{Stage},B$. Set the admitted
$\mathsf{Stage}$ and $B$ store fibres to have $d=1$; the $A$ fibre
permits either bit. This is a definition of the well-formed state space,
not an inference about arbitrary semantic states from reachability.
The authority tokens have no nontrivial private state.

All operations below are finite Dirac method tables with a public unary
node, scalar cost one, and the displayed actual-token transition. Public
labels include the method name and its returned value. Every unaffected
resource explicitly carried by a table remains in its declared output;
there is no generic implicit capability frame.
\begin{description}
\item[Stage.]
$(r:A_h,k:\mathsf{Fresh})\to(r:\mathsf{Stage}_h,k:\mathsf{Armed})$,
with $(A,x,d)\mapsto(\mathsf{Stage},1-x,1)$ and result
$\mathsf{staged}$.
\item[Rollback.]
$(r:\mathsf{Stage}_h,k:\mathsf{Armed})\to r:A_h$, consuming $k$,
with $(\mathsf{Stage},y,1)\mapsto(A,1-y,1)$ and result
$\mathsf{rolledBack}$.
\item[Commit.]
$(r:\mathsf{Stage}_h,k:\mathsf{Armed})\to r:B_h$, consuming $k$,
retaining the encoded payload and audit bit, with result
$\mathsf{committed}$.
\item[Cancel.]
$(r:A_h,k:\mathsf{Fresh})\to r:A_h$, consuming $k$, leaving the
service store unchanged, with result $\mathsf{cancelled}$. This is the
actual two-token disposal table, not an untyped omission of $k$.
\item[Read.]
For $h>0$, $r:A_h\to(\mathsf{Bool},r:A_{h-1})$, returning $x$ and
leaving $(A,x,d)$ unchanged.
\item[Write a Boolean $b$.]
For $h>0$, $r:A_h\to r:A_{h-1}$, changing $(A,x,d)$ to $(A,b,d)$
and returning $\mathsf{written}$.
\item[Health.]
For $h>0$, $r:A_h\to r:A_{h-1}$, leaving the store unchanged and
returning $\mathsf{AReady}$.
\item[Audit.]
$(r:A_h,a:\mathsf{Audit1})\to(\mathsf{Bool},r:A_h,a:\mathsf{Audit0})$,
returning $d$ and leaving the service store unchanged.
\end{description}
Audit remains in the global signature. No method creates one of these
identities, restores a consumed $k$, grants $\mathsf{Audit1}$, decreases
$d$, or turns $B$ back into $\mathsf{Stage}$. Read/write analogues in
$\mathsf{Stage}$ or $B$ are not needed and are not supplied in this
small extension. They cannot be assumed as extra continuations without
checking a new signature.

The exact recovery domain is
\begin{equation}\label{eq:recovery-domain}
 D_{\mathrm{rec}}=\{(r:\mathsf{Stage}_h,k:\mathsf{Armed},
                  (\mathsf{Stage},y,1)):0\le h\le N, y\in\{0,1\}\},
\end{equation}
with any explicitly admitted disjoint carried resources. Rollback is
total on this typed finite domain. A committed $B$ state, a receipt
without owned $k$, or an already consumed $k$ is outside it. A different
API could declare a refusal result, but no refusal is silently identified
with successful recovery here.

\subsection{The right baseline and the residual relation}
Start at $r:A_h$ with state $(A,x,d)$ and owned $k:\mathsf{Fresh}$.
Let $D$ be the pure token-preserving map that discards an ordinary
receipt and returns public unit. Use the explicit boundary prefixes
$\mathsf{stage};D;\mathsf{rollback};D$ and
$\mathsf{cancel};D$. The first $D$ supplies rollback's unit-data input;
the final maps remove the unequal $\mathsf{rolledBack}$ and
$\mathsf{cancelled}$ receipt values. They do not erase public event
nodes or costs. With equal otherwise retained public memory $\Gamma$,
these prefixes lead to
\begin{align}\label{eq:residual-states}
 \mathsf{stage};D;\mathsf{rollback};D:&\quad
 s_{\mathrm{rec}}=(r:A_h,(A,x,1),\text{no }k),\notag\\
 \mathsf{cancel};D:&\quad
 s_{\mathrm{base}}=(r:A_h,(A,x,d),\text{no }k).
\end{align}
The boundary comparison requires identical public unit data and equal
retained public $\Gamma$; it does not cover an observer that carries a
different pre-boundary receipt or event history in its public memory.
Such an observer can already distinguish the prefixes. The baseline
explicitly disposes of unused recovery authority. Comparing
instead to a state still owning $k:\mathsf{Fresh}$ would change the
future interface and invalidate the claimed equality. The request
counter $h$ is unchanged because neither prefix uses a service request;
rollback does not replenish previously spent request credits.

At the residual service kind define
\begin{equation}\label{eq:residual-relation}
 (A,x,d_0)\ \mathcal R_{r,A_h}\ (A,x',d_1)
       \quad\Longleftrightarrow\quad x=x'.
\end{equation}
For the functional regime, take only the resource-indexed interfaces
reachable from the stipulated residual ownership without $k$ or an
$\mathsf{Audit1}$ token; that family is closed under the declared
transitions and disjoint restrictions. All public configurations,
actual ownership, other private entries and remaining counters must agree.
Equivalently extend this relation to the full global signature with empty
relation fibres on $k:\mathsf{Fresh}$, $k:\mathsf{Armed}$,
$a:\mathsf{Audit1}$ and service phases $\mathsf{Stage},B$; retain the
usual equality relation on an inert $a:\mathsf{Audit0}$ if present.
Thus unavailable methods have no related input, rather than an omitted
local obligation. This relation forgets precisely the
protected audit bit, not the functional payload. The states
in~\eqref{eq:residual-states} are related for every admitted $x,d,h$.

The prefixes are not equivalent observations. Staging and rollback emit
two public nodes and cost $2$; cancel emits a different single node and
costs $1$. The residual theorem begins \emph{after} this boundary and
does not erase that public history. With paired functional continuation
cost $c$, the total costs are exactly $2+c$ and $1+c$. This is overhead
one relative to the cancel baseline, not a full public-tree refinement,
a cost refund, or equality with doing nothing.

\subsection{All functional continuations preserve the residual relation}
Let $\mathcal C_{\mathrm{fun}}(h)$ be all finite O4 closing contexts
starting at the residual public interface $r:A_h$ without owned
$a:\mathsf{Audit1}$ and with no operation able to acquire such authority.
They retain all other well-typed operations of the fixed signature.
Since $k$ is absent and cannot be created, staging, commit, rollback and
cancel cannot be invoked at this boundary. Since the audit authority is
absent, audit is not typable. This is a capability-defined class, not a
post hoc deletion of a public test from an otherwise unrestricted
interface. A different initial interface with audit authority is
explicitly analyzed below.

Contexts can read, write and test health repeatedly within the actual
counter, store answers and code aliases, adapt their next operation or
stop decision, sample client randomness, run declared unrelated work in
disjoint batches, and invoke every permitted ranked template. Observations
retain the full public causal tree, public final leaf data and owned token
kinds, and exact continuation cost. No private management-event erasure
is needed for this example; all displayed method events are public.

\begin{theorem}[Residual continuation equivalence]\label{thm:recovery}
For every $C\in\mathcal C_{\mathrm{fun}}(h)$ and every joint initial
coupling supported on the expanded relation $\mathcal R_E$, the two
observed continuation laws are equal. In particular,
\begin{equation}\label{eq:recovery-equivalence}
 \Obs(C[s_{\mathrm{base}}])=\Obs(C[s_{\mathrm{rec}}]).
\end{equation}
The statement is conditional on successful authorized rollback in
$D_{\mathrm{rec}}$ and on the stated residual capabilities.
\end{theorem}
\begin{proof}
We verify every available local method at arbitrary related stores.
Read returns the same $x$, preserves each side's possibly different
$d$, and decrements the same counter. Its response-labelled event and
cost agree. Write receives the same public Boolean $b$, installs that
same logical value, preserves both audit bits and decrements the same
counter; its event and cost agree. Health returns the same readiness
value, preserves the stores and decrements the same counter, again with
equal event and cost. Thus each Dirac pair is an exact-cost local
matrix with $\mathcal R$-related leaves. Other available primitive tables
satisfy L3 by their declared independence from $r$'s audit entry and their
explicit carried-token behavior. Audit has no well-typed input in this
class; neither does an operation needing the absent $k$.

The pure and ranked-code clauses of L4 are unchanged. The given initial
joint coupling is L1, including correlations with other tokens or public
data. The exact-cost form of Theorem~\ref{thm:replacement} therefore
applies to every such context. Its native paired-history realization has
both projection legs cost-exact. Final private-store projection identifies
the related leaves and yields~\eqref{eq:recovery-equivalence}.
\end{proof}
This class has meaningful tests. If $h\ge1$, read distinguishes residual
payload $0$ from payload $1$ with probability one. If $h\ge2$,
write$(1)$ followed by read checks a persistent update. A read followed
by a data-dependent write tests adaptive interaction and resource change.
The theorem preserves these public results and exact costs; it does not
obtain equivalence by removing every state-sensitive continuation.

\subsection{The authorized audit distinguisher}
If the residual interface also owns $a:\mathsf{Audit1}$, take initial
$d=0$ and the actual typed continuation
\[
 \mathsf{audit}(r,a);\quad\Ret(\text{returned Boolean})
 \quad [r:A_h,\ a:\mathsf{Audit0}\text{ retained}].
\]
It returns $0$ from $s_{\mathrm{base}}$ and $1$ from
$s_{\mathrm{rec}}$, each with probability one, the same cost and the same
resource transitions. Thus the broader residual equivalence is false.
The local audit matrix cannot satisfy equality of public replies; the
failure is identified in the premise rather than hidden in the proof.
When the initial audit bit is already $1$, these particular residual
stores coincide, although their prefixes still differ in events and cost.

Equivalently, the projection $(A,x,d)\mapsto(A,x)$ preserves the
functional method tables but is not a homomorphism for the full
audit-capable signature: the public audit reply cannot be computed from
that quotient. This is a concrete boundary on early state erasure and
on which continuations a residual relation supports.

\subsection{Why evidence and model revision are not recovery}
The result $\mathsf{staged}$ is an ordinary receipt. A proof that
rollback is possible, a model simulation, or a posterior assigning high
probability to recovery is likewise not an owned $k:\mathsf{Armed}$.
Only the actual rollback table with its typed authority premise provides
the state transition. This follows the frozen separation of evidence,
code and capability; there is no inverse-execution operation generated
by a theorem or a prediction.

Suppose an earlier model falsely said rollback removed every durable
effect, and later observation establishes that $d$ remains set. Correcting
the model is a revision of the transition table. Audit-authorized clients
distinguish the old and corrected residual laws, so the correction is
not automatically exact-law refinement in either direction. The restricted
functional equivalence above has its own proof; it is not a consequence
of naming the correction learning. Recovery is an application property,
not a definition imposed on every semantic world model.

\section{Existing theory and the exact increment}\label{sec:prior}
Probabilistic couplings as relational proof witnesses are established:
Barthe and collaborators connect them to probabilistic relational Hoare
logic~\cite{Coupling}. Contextual reasoning for higher-order probability
and state is also much broader than this finite fragment. Bizjak and
Birkedal construct step-indexed logical relations for a language with
polymorphism and recursive types, and extend their treatment to dynamic
first-order references~\cite{Bizjak}. Relators and compositional
applicative reasoning for monadic effects are established categorical
techniques~\cite{Relators}; using a relational lifting is not itself a
novel contribution.

Approxis supplies higher-order approximate relational separation logic
with state and exact contextual-equivalence applications, mechanized in
Coq/Iris~\cite{Approxis}. The stronger current baseline Foxtrot proves
contextual refinement for higher-order concurrent probabilistic programs
with higher-order local state; its results are mechanized in Rocq/Iris
\cite{Foxtrot}. O5 has no claim to supersede that expressiveness or
verification strength. In particular, finite ranked callbacks are not
a substitute for Foxtrot's concurrency or higher-order store results.

Cost-sensitive relational reasoning is likewise established in RelCost
\cite{RelCost}; Tachis provides mechanized higher-order probabilistic
reasoning about expected and amortized costs~\cite{Tachis}. The present
supportwise cost order inside a joint coupling is a different specified
guarantee, not an assertion that expected-cost theory is inadequate or
that generic cost refinement is new.

The increment here is a self-contained specialization to the contracted
language and its retained carriers: explicit capability-dependent local
matrices, full open constructor closure, the paired-history realization
with one common reference cost and actual O1 endpoint equations, its
O3 representation comparison, and a sharp residual recovery example.
The result claims no publication novelty for coupling, simulation,
logical relations, contextual equivalence, relational lifting or
compensation. Primary-source pages for the cited external baselines were
checked on 1 November 2026; this is not an exhaustive literature survey.

\section{Acceptance ledger and remaining boundaries}\label{sec:acceptance}
\begin{longtable}{@{}p{0.29\textwidth}p{0.66\textwidth}@{}}
\toprule
O5 obligation&Established result\\
\midrule
\endhead
Context class before replacement&Section~\ref{sec:regime}; all finite
well-typed O4 contexts at fixed interface, signature and ranked library,
including stored aliases, repeated private randomness, adaptive resources
and disjoint batches.\\
Exact observation and order&Equations~\eqref{eq:observation} and
\eqref{eq:obsorder}; public causal trees and capability leaves match,
external costs are ordered, private state is projected only at closing.\\
Noncircular local premise&L1--L4 and~\eqref{eq:local-matrix}; finite
method/input matrices and state-relation checks, not a hypothesis about
all clients.\\
Quantified open theorem&Theorem~\ref{thm:replacement}; full structural
and rank induction, actual dependent leaves and arbitrary joint input
laws. No finite client enumeration replaces it.\\
Categorical connection&Lemma~\ref{lem:raw-lift} and
Section~\ref{sec:cells}; native endpoint equations over relation carriers,
reference-exact and candidate-lowering costs, actual O1 occurrence
composition, scoped law-level vertical gluing.\\
Representation comparison&Theorem~\ref{thm:representation}; explicit
finite pushforward into the O3 operational/law/QBS execution models,
without premature private-state erasure.\\
Nontrivial positive pair&Corollary~\ref{cor:persistent}; distinct private
carriers and raw initialization histories, persistent reused randomness,
strictly lower per-call cost.\\
Matching coarse endpoints fail&Section~\ref{sec:examples}; shared and
fresh services have identical one-call marginals and resource endpoints
but different two-call equality probabilities.\\
Recovery domain and authority&Equation~\eqref{eq:recovery-domain}; an
actual staged state and owned armed token, no receipt-to-reset rule.\\
Residual guarantee and boundary&Theorem~\ref{thm:recovery}; all functional
continuations after baseline authority disposal; authorized audit
distinguishes; prefix events and exact overhead remain.\\
Zero-probability scope&Section~\ref{sec:zeros}; law guarantees do not
imply all-path budgets; candidate projection coverage is the explicit
additional sufficient premise.\\
\bottomrule
\end{longtable}

The private-state relation is a sufficient tokenwise invariant, not every
possible relational invariant. Costs cannot secretly influence public
control. Public code remains inspectable, and every imported body must
obey the actual declared interface and finite call discipline. Continuous
state, arbitrary QBS measurable coupling selection, dynamic name
allocation, unrestricted recursion, unknown foreign callback semantics,
physical interleaving, fairness, full abstraction and a general new
cell-coherence theorem remain outside scope. None is used implicitly to
obtain the finite result.

\section{Source checks and verification status}\label{sec:checks}
The accompanying archive contains this uncompiled TeX source,
standard-library exact finite regression programs, their recorded output,
a source-only checker, a reproducible explicit-inventory archive builder,
a README, a manifest and checksums. The checks exercise local simulation
marginals and support, the raw conditional lift and native cost equations,
conditional products with correlated inputs, the sharing distinguisher,
zero-weight scope and the recovery transition/invariance tables.
They supplement the proofs; they do not enumerate every admitted context
or decide the arbitrary-real-weight theorem.

Source checks cover balanced braces and environments, internal references,
bibliography keys, principal statement labels, and unresolved drafting
markers. Archive checks verify inventory and file hashes. No TeX engine
is invoked. There is therefore no compiled PDF, rendered-layout approval,
proof-assistant certification, physical service execution or externally
tested recovery claim. The uncompiled source is the requested artifact.

\begin{thebibliography}{99}
\bibitem{Contract}
\emph{Expanded Research Contract Version 0.1}. Project source,
1 November 2026, O5, lines 529--570.
\texttt{Expanded\_Research\_Contract\_v0\_1.tex}.

\bibitem{O1}
\emph{General Independent Interaction Criterion}. Project proof dossier,
1 November 2026. Augmented families, native endpoint cells and the mixed
occurrence construction.
\texttt{General\_Independent\_Interaction\_Criterion.tex}.

\bibitem{O3}
\emph{Plural Semantics for Result Dependent Protocols}. Frozen project
proof dossier, 1 November 2026. CF-0.1 and its retained, operational,
finite-law and QBS comparisons.
\texttt{Plural\_Semantics\_for\_Result\_Dependent\_Protocols.tex}.

\bibitem{O4}
\emph{Typed Protocols and Bounded Higher Order Interaction}. Frozen
project proof dossier, 1 November 2026. Finite dependent typing,
substitution, private stores and ranked code.
\texttt{Typed\_Protocols\_and\_Bounded\_Higher\_Order\_Interaction.tex}.

\bibitem{Coupling}
G. Barthe, T. Espitau, B. Gr\'{e}goire, J. Hsu, L. Stefanesco and
P.-Y. Strub. \emph{Relational reasoning via probabilistic coupling}.
arXiv:1509.03476v2, 2015. \url{https://arxiv.org/abs/1509.03476}.

\bibitem{Bizjak}
A. Bizjak and L. Birkedal. \emph{Step-Indexed Logical Relations for
Probability}. FoSSaCS 2015; extended version arXiv:1501.02623v2.
\url{https://arxiv.org/abs/1501.02623}.

\bibitem{Relators}
U. Dal Lago, F. Gavazzo and P. B. Levy. \emph{Effectful Applicative
Bisimilarity: Monads, Relators, and Howe's Method}. LICS 2017;
arXiv:1704.04647. \url{https://arxiv.org/abs/1704.04647}.

\bibitem{Approxis}
P. G. Haselwarter, K. H. Li, A. Aguirre, S. O. Gregersen, J. Tassarotti
and L. Birkedal. \emph{Approximate Relational Reasoning for Higher-Order
Probabilistic Programs}. PACMPL 9, POPL, article 41, January 2025.
DOI: 10.1145/3704877. \url{https://arxiv.org/abs/2407.14107}.

\bibitem{Foxtrot}
K. H. Li, A. Aguirre, J. Tassarotti and L. Birkedal.
\emph{Contextual Refinement of Higher-Order Concurrent Probabilistic
Programs}. PACMPL 10, PLDI, article 187, June 2026.
DOI: 10.1145/3808265.
\url{https://iris-project.org/pdfs/2026-pldi-foxtrot.pdf}.

\bibitem{RelCost}
E. \c{C}i\c{c}ek, G. Barthe, M. Gaboardi, D. Garg and J. Hoffmann.
\emph{Relational Cost Analysis}. POPL 2017.
\url{https://people.mpi-sws.org/~dg/papers/popl17.pdf}.

\bibitem{Tachis}
P. G. Haselwarter, K. H. Li, M. de Medeiros, S. O. Gregersen,
A. Aguirre, J. Tassarotti and L. Birkedal.
\emph{Tachis: Higher-Order Separation Logic with Credits for Expected
Costs}. PACMPL 8, OOPSLA2, article 313, October 2024.
DOI: 10.1145/3689753.
\url{https://iris-project.org/pdfs/2024-oopsla-tachis.pdf}.
\end{thebibliography}
\end{document}
